\documentclass[12pt]{article}

\usepackage[spanish]{babel}
\usepackage[utf8]{inputenc}
\usepackage[T1]{fontenc}
\usepackage{geometry}
\usepackage{amsmath}
\usepackage{listings}
\usepackage{xcolor}
\usepackage{hyperref}
\usepackage{array}
\usepackage{longtable}
\usepackage{tcolorbox}

\geometry{margin=2.5cm}

% Estilo código C++
\lstset{
    language=C++,
    basicstyle=\ttfamily\small,
    keywordstyle=\color{blue},
    commentstyle=\color{gray},
    stringstyle=\color{green!40!black},
    numbers=left,
    numberstyle=\tiny,
    breaklines=true,
    frame=single,
    showstringspaces=false
}

% Caja para respuestas
\newtcolorbox{respuesta}{
colback=gray!5,
colframe=black!40,
title=Respuesta,
}

\begin{document}

% ---------------- PORTADA ----------------

\begin{titlepage}

\centering

\vspace*{2cm}

{\Large \textbf{UNIVERSIDAD EAFIT}}

\vspace{0.5cm}

{\large Ingeniería de Sistemas}

\vspace{0.5cm}

{\large Análisis y Diseño de Algoritmos}

\vspace{2cm}

{\Huge \textbf{PRIMER TRABAJO PRÁCTICO}}

\vspace{0.5cm}

{\Large Fuerza Bruta y Backtracking en C++}

\vspace{3cm}

\begin{flushleft}

\textbf{Integrantes:} Sebastian Salazar, Andres Velez y Nathalia Cardoza. \\

\vspace{0.3cm}

\textbf{Profesor:} Carlos Alberto Alvarez Henao. \\

\vspace{0.3cm}

\textbf{Fecha:} 5 de Abril de 2026 \\

\end{flushleft}

\vfill

\end{titlepage}

\tableofcontents

\newpage

% ---------------- INTRO ----------------

\section{Introducción}

En este trabajo se analizan e implementan dos problemas clásicos utilizando los enfoques de \textbf{fuerza bruta} y \textbf{backtracking}. El objetivo principal es comprender las diferencias entre ambos métodos en términos de eficiencia, complejidad y aplicabilidad en problemas combinatorios.

El primer ejercicio consiste en generar permutaciones con restricciones, mientras que el segundo aborda el problema de coloración de grafos con $k$ colores.

El desarrollo incluye pseudocódigo, implementación en C++, análisis de complejidad y comparación experimental entre ambos enfoques.

% ---------------- TEORIA ----------------

\section{Conceptos Teóricos}

\subsection{Fuerza Bruta}

La fuerza bruta consiste en evaluar todas las posibles soluciones hasta encontrar aquellas que cumplen las condiciones del problema.

Características:

\begin{itemize}
\item Explora todo el espacio de búsqueda
\item Garantiza encontrar la solución
\item Puede ser costoso computacionalmente
\end{itemize}

\subsection{Backtracking}

El backtracking mejora la fuerza bruta al descartar soluciones inválidas de manera anticipada.

Características:

\begin{itemize}
\item Utiliza recursividad
\item Reduce el número de combinaciones exploradas
\item Aplica poda en el árbol de búsqueda
\end{itemize}

\subsection{Diferencias}

\begin{center}
\begin{tabular}{|c|c|c|}
\hline
Método & Explora todo & Usa poda \\
\hline
Fuerza bruta & Sí & No \\
Backtracking & No & Sí \\
\hline
\end{tabular}
\end{center}

% ---------------- EJERCICIO 1 ----------------

\section{Ejercicio 1: Permutaciones con Restricción}

\subsection{Descripción}

Dado un conjunto de enteros distintos, se deben generar todas las permutaciones posibles que cumplan:

\begin{equation}
P[i] \leq 2P[i+1]
\end{equation}

\subsection{Idea del algoritmo}

Explicar con tus palabras cómo funciona.

\begin{respuesta}

Escribir explicación aquí.

\end{respuesta}

\subsection{Por qué es fuerza bruta}

\begin{respuesta}

Explicar por qué evalúa todas las permutaciones.

\end{respuesta}

\subsection{Pseudocódigo}

\begin{verbatim}
leer n
leer arreglo

ordenar arreglo

hacer todas las permutaciones

verificar restriccion

contar validas
\end{verbatim}

\subsection{Implementación en C++}

\begin{lstlisting}
// pegar codigo ejercicio 1
\end{lstlisting}

\subsection{Ejemplo}

Entrada:

\begin{verbatim}
3
1 3 4
\end{verbatim}

Salida esperada:

\begin{verbatim}
Permutaciones validas:
...

Total generadas = 6
Total validas = 3
\end{verbatim}

\subsection{Complejidad}

\begin{respuesta}

Tiempo:

O(n!)

porque se generan todas las permutaciones posibles.

Espacio:

O(n)

\end{respuesta}

\subsection{Tabla de tiempos}

\begin{center}
\begin{tabular}{|c|c|}
\hline
n & tiempo (ms) \\
\hline
6 & \\
8 & \\
10 & \\
\hline
\end{tabular}
\end{center}

\subsection{Preguntas guía}

\textbf{¿Qué ocurre con la proporción de permutaciones válidas cuando crece n?}

\begin{respuesta}

Escribir respuesta.

\end{respuesta}

\textbf{¿Dónde se podría aplicar poda con backtracking?}

\begin{respuesta}

Escribir respuesta.

\end{respuesta}

\textbf{¿Por qué next\_permutation genera n! permutaciones?}

\begin{respuesta}

Escribir respuesta.

\end{respuesta}

% ---------------- EJERCICIO 2 ----------------

\section{Ejercicio 2: Coloración de Grafos}

\subsection{Descripción}

Se desea asignar colores a los vértices de un grafo sin que vértices adyacentes tengan el mismo color.

\subsection{Idea del algoritmo}

\begin{respuesta}

Explicar lógica del backtracking.

\end{respuesta}

\subsection{Pseudocódigo}

\begin{verbatim}
asignar color a cada nodo

si hay conflicto:
    probar otro color

si no hay colores posibles:
    volver atras
\end{verbatim}

\subsection{Implementación Backtracking}

\begin{lstlisting}
// pegar codigo backtracking
\end{lstlisting}

\subsection{Implementación Fuerza Bruta}

\begin{lstlisting}
// pegar codigo fuerza bruta
\end{lstlisting}

\subsection{Ejemplo}

Entrada:

\begin{verbatim}
n = 4
k = 3
\end{verbatim}

Salida:

\begin{verbatim}
Total soluciones = 18
\end{verbatim}

\subsection{Comparación}

\begin{center}
\begin{tabular}{|c|c|}
\hline
Método & Complejidad \\
\hline
Fuerza bruta & $O(k^n)$ \\
Backtracking & menor o igual a $O(k^n)$ \\
\hline
\end{tabular}
\end{center}

\subsection{Tabla de tiempos}

\begin{center}
\begin{tabular}{|c|c|c|}
\hline
grafo & fuerza bruta & backtracking \\
\hline
1 & & \\
2 & & \\
\hline
\end{tabular}
\end{center}

\subsection{Complejidad}

\begin{respuesta}

Fuerza bruta:

O(k^n)

Backtracking:

explora menos combinaciones.

Espacio:

O(n)

\end{respuesta}

\subsection{Preguntas guía}

\textbf{¿Cuántos nodos explora cada método?}

\begin{respuesta}

Responder aquí.

\end{respuesta}

\textbf{¿Qué grafos benefician más al backtracking?}

\begin{respuesta}

Responder aquí.

\end{respuesta}

\textbf{¿Cómo saber solo si existe solución?}

\begin{respuesta}

Responder aquí.

\end{respuesta}

\textbf{Relación con número cromático}

\begin{respuesta}

Responder aquí.

\end{respuesta}

% ---------------- CONCLUSIONES ----------------

\section{Conclusiones}

\begin{respuesta}

Escribir conclusiones:

• diferencia entre fuerza bruta y backtracking

• eficiencia

• aprendizaje obtenido

\end{respuesta}

% ---------------- REFERENCIAS ----------------

\section{Referencias}

\begin{itemize}

\item Material de clase

\item Documentación C++

\item Libro de algoritmos

\item ChatGPT (si decides citarlo)

\end{itemize}

\end{document}